%!TEX root = ./main.tex
%
The key role in our platform is played by \emph{service brokers} that
are used by \emph{service providers} to register their services, and by 
\emph{service clients} to dynamically \emph{discover} and \emph{bind} to 
services satisfying their requirements.
%
These roles are typically advocated in SOAs; the originality of
our platform however is the use of \emph{behavioural contracts} to
support \emph{semantic-based} composition of services.
%
More precisely, the composition mechanism realised by service
brokers of \SEArch relies on a notion of behavioural compliance of
contracts which we briefly review here.

To support SOC, our platform offers a mechanism for populating
service repositories and composing service-based applications.
%
The registration of a service (see \cref{fig:service-registering-procedure}) 
is conceptually simple: the service provider sends the contract and the 
unique resource identifier of the provided service, to broker, through its 
local \SEArch middleware.
%
\begin{figure}[t!]
  % First subfigure (left)
  \begin{subfigure}{\textwidth}
  \centering
  \includegraphics[width=0.65\textwidth]{registering}
    \caption{\label{fig:service-registering-procedure} Service registration procedure}
  \end{subfigure}
  \\[2em]
  \begin{subfigure}{\textwidth}
  \centering
  \includegraphics[width=0.85\textwidth]{brokering}
    \caption{\label{fig:service-brokering-procedure} Service brokering procedure}
  \end{subfigure}
  \caption{ \label{fig:execution} Service management procedures in \SEArch}
\end{figure}
%\begin{figure}[t]
%  \centering
%  \tikzset{
%	 ser/.style = {
%		rectangle,
%		draw=blue!50,
%		opacity=.7,
%		line width=.5pt,
%		align=center,
%		fill=white,
%		minimum width = 2cm,
%		minimum height = 1cm, 
%		rounded corners=0.1cm,
%		inner sep=10pt
%	 },
%	 mid/.style = {
%		rectangle,
%		draw,
%		opacity=.7,
%		line width=.5pt,
%		align=center,
%		fill=yellow!20,
%		minimum size=4mm,
%		inner sep=2pt
%	 },
%	 arr/.style = {
%		->, >=latex, dashed
%	 }
%  }
%  \begin{tikzpicture}
%	 \node (client) [ser, label={above:Service Client}]{};
%	 \node (cmid) [mid] at (client.center){Middleware};
%	 \node at (cmid.east){\tiny$\bullet$};
%	 %
%	 \node (provider) [ser,label={above:Service Provider},right = 3cm of client]{};
%	 \node (pmid) [mid] at (provider.center){Middleware};
%	 \node at (pmid.east){\tiny$\bullet$};
%	 \node at (pmid.west){\tiny$\bullet$};
%	 %
%	 \node (broker) [ser, label={right:Service Broker}, below right = of client, xshift = -3cm, minimum width=8cm]{};
%	 \node (bmid) [rectangle, draw, label={below:Compliance}, yshift = .01cm] at ($(broker.center) + (-2.5cm,0)$){$\mathsf{Pr} \vdash \mathsf{R}$};
%	 \node (repo) [rectangle, draw, label={below:Repository}, yshift = .01cm] at ($(broker.center) + (2cm,0)$){\tiny$\conf{\mathsf{R}_1,u_1}\cdots\conf{\mathsf{R}_n,u_n}$};
%	 %
%	 \path[draw] ($(pmid.east) + (1cm, 0)$) -- (pmid.east);
%	 \path[draw] (pmid.west) -- (cmid.east);
%	 \path[arr] ($(bmid.east) + (0,.1cm)$) edge node[above]{query $\mathsf{R}$} ($(repo.west) + (0,.1cm)$);
%	 \path[arr] ($(repo.west) + (0,-.1cm)$) edge node[below]{$\conf{\mathsf{Pr},u}$} ($(bmid.east) + (0,-.1cm)$);
%	 \path[arr] ($(broker.north west) + (.5cm,0)$) edge node[left]{$\conf{\mathsf{Pr},u}$} ($(cmid.south west) + (.5cm,0)$);
%	 \path[arr] ($(cmid.south west) + (1cm,0)$) edge node[right]{$\mathsf{R}$} ($(broker.north west) + (1cm,0)$);
%	 \path[arr] ($(broker.north east) + (-.5cm,0)$) edge node[right]{$\conf{\mathsf{Pr},u}$} ($(pmid.south east) + (-.5cm,0)$);
%	 \path[arr] ($(pmid.south east) + (-1cm,0)$) edge node[left]{$\mathsf{R}$} ($(broker.north east) + (-1cm,0)$);
%  \end{tikzpicture}
%  \caption{Service execution procedure in \SEArch}
%  \label{fig:execution}
%\end{figure}
%
The execution process is significantly more
complex since its underlying computational model are
\emph{Asynchronous Relational Networks}~(ARNs)~\cite{fiadeiro:tcs-503}.
%
This theory envisages \emph{processes} and \emph{communication
  channels} exposing \emph{ports} with behavioural contracts:
processes' ports, called \emph{provides-points}, can be dynamically
\quo{coalesed} with ports of communication channels, called
\emph{requires-points}, provided that the behavioural contracts
associated with merged ports are \quo{semantically compliant}.
%
The notion of compliance adopted is defined in terms of
contracts formalised in terms of a bisimilarity relation
between \emph{extended Communicating Finite State Machines} (extended
CFSM for short)~\cite{lopezpombo:coordination25}, a variant of
\emph{communicating finite-state machines}~\cite{brand:jacm-30_2}
enabling the specification of functional and non-functional behaviour.
%
We now describe our execution process.

When launched on a \SEArch node, a service registers the extended
CFSMs of its exposed requirement contracts on the local
\emph{middleware}, which will mediate the communication with other
services.
%
Basically, the middleware, intercepts the communications a
local service attempts over a communication channel, establishes a
communication session for that channel (in case the session does not 
yet exists), and triggers the dynamic reconfiguration sketched in
\cref{fig:service-brokering-procedure} which we now describe:
%
\begin{enumerate}
  \newcommand{\blt}[1]{\raisebox{3pt}{\tikz{\node[circle,fill=black,text=white,inner sep=1pt, scale=.9]{#1}}}}
\item[\blt 1] the middleware sends the service broker the contract
  $\{\mathsf{R_1}, \ldots, \mathsf{R_m}\}$ of the communication
  channel, consisting of one extended CFSM for each participant
  required by the service client;
\item[\blt 2] for each extended CFSM $\mathsf{R_i}$ in the contract, the
  service broker queries the service repository for candidates;
\item[\blt 3] the service repository returns a set of candidates\footnote{%
	 The query is resolved by filtering the extended CFSMs in the
	 repository according to heuristics based on structure of
	 $\mathsf{R_i}$.
	 %
  }
  $\{\conf{\mathsf{P_{i, 1}}, \mathsf{u_{i, 1}}}, \ldots,
  \conf{\mathsf{P_{i, k_i}}, \mathsf{u_{i, k_i}}}\}$ where, for all
  $1 \leq j \leq k_i$, $\mathsf{P_{i, j}}$ is an extended CFSM and
  $\mathsf{u_{i, j}}$ is the URI of the service;
\item[\blt 4] the broker checks whether there is a provision contract $\conf{\mathsf{P_{i, j}}, \mathsf{u_{i, j}}}$, with $1 \leq j \leq k$,
  compatible with the requirement contract $\mathsf{R_i}$, for all $1 \leq i \leq m$; this is
  done by resorting to a bisimilarity check; once the service broker has
  found services satisfying all the requirement contracts, it returns
  the set of URIs $\{\mathsf{u_1}, \ldots, \mathsf{u_m}\}$ to the middleware;
\item[\blt 5] finally, the middleware opens a communication with the
  middleware of each service returned by the provider requiring the
  execution of the corresponding service.
\end{enumerate}
%
Once the composition is achieved, the middleware interacts with the
middlewares of the partner services; and the execution process
proceeds.
%
Notice that during the execution new requirements might crop up; for
instance, because some brokered services need further
services.\footnote{%
  This implies that a service may have one or more required-points
  associated to communication channels of its own.
  %
}
% 
This will initiate a new brokerage phase to discover the newly
required services.

The schematic view discussed before establishes several requisites
over the implementation of the middleware and the service broker which
are omitted as inessential here; interested readers are referred
to~\cite{lopezpombo:coordination24} for details.


%%% Local Variables:
%%% mode: latex
%%% TeX-master: "main"
%%% End:
